\chapter{Область действия электромагнитного зарядового тождества}
\label{ch:v8-baryon-em-total-charge-identity-scope-gate}

Проверим отдельно алгебраическое тождество, его превращение в энергию и
знак нейтрон-протонного вклада. Это три разных утверждения, которые нельзя
сливать в одну физическую теорему.

\section{Тождество на полном трёхкварковом носителе}

На одном слабом дублете оператор заряда равен
\begin{equation}
 Q=\operatorname{diag}\!\left(\frac23,-\frac13\right).
 \label{eq:v8-baryon-em-one-particle-charge}
\end{equation}
На $(\mathbb C^6)^{\otimes3}$ определим
\begin{equation}
 A=\sum_{r=1}^{3}(Q^{(r)})^2,
 \qquad
 C=\sum_{r\ne s}Q^{(r)}Q^{(s)},
 \qquad
 Q_{\rm tot}=\sum_{r=1}^{3}Q^{(r)}.
 \label{eq:v8-baryon-em-self-pair-total}
\end{equation}
Тогда непосредственно
\begin{equation}
 A+C=Q_{\rm tot}^2.
 \label{eq:v8-baryon-em-total-charge-identity}
\end{equation}
Это операторное тождество на всех $216$ базисных состояниях; цветовая
эпсилон-проекция для него несущественна. Если $n_d$ --- число нижних
компонент дублета, то
\begin{equation}
 Q_{\rm tot}=2-n_d,
 \qquad
 A=\frac{4-n_d}{3},
 \qquad
 C=(2-n_d)^2-\frac{4-n_d}{3}.
 \label{eq:v8-baryon-em-epsilon-pattern}
\end{equation}
Поэтому последовательность $Q_{\rm tot}^2$ при $n_d=0,1,2,3$ равна
$(4,1,0,1)$ точно.

\begin{proposition}[Зарядовое тождество]
Формула $A+C=Q_{\rm tot}^2$ и рисунок $(4,1,0,1)$ являются строгими
алгебраическими утверждениями данного зарядового представления.
\end{proposition}

\section{Следовая нормировка}

Канонический генератор $Q=T_3+Y$ имеет на каждом из двух углов полный след
квадрата $14/3$. Для КМС-плотности
$\rho=aI_{11}\oplus bI_{10}$ получаем
\begin{equation}
 T=\Tr(\rho Q^2)=\frac{14}{3}(a+b)
 =\frac{14(1+x)}{3(11+10x)}.
 \label{eq:v8-baryon-em-common-trace-norm}
\end{equation}
Это точная безразмерная норма направления заряда. Она не является сама по
себе электромагнитной энергией и не задаёт переход к физическим единицам.

\section{Общий перестановочно-инвариантный электростатический оператор}

Даже после запрета различать три копии наиболее общий оператор из данных
двух зарядовых структур имеет два коэффициента:
\begin{equation}
 H_{\mu,\lambda}=\frac1T(\mu A+\lambda C).
 \label{eq:v8-baryon-em-two-coefficient-form}
\end{equation}
Равенство $H_{\mu,\lambda}=Q_{\rm tot}^2/T$ на всём носителе требует и
достаточно требует
\begin{equation}
 \mu=1,
 \qquad
 \lambda=1.
 \label{eq:v8-baryon-em-collapse-condition}
\end{equation}
Ни зарядовое представление, ни эпсилон-синглет, ни общий след не выводят
это равенство коэффициентов. Пространственные коэффициенты парного ядра,
дипольные и магнитные операторы в проверяемой формуле отсутствуют по
определению, а не вследствие отдельного запрета.

\begin{nogocriterion}[Граница зарядового тождества]
Алгебраическое равенство $A+C=Q_{\rm tot}^2$ нельзя повышать до
универсальной электромагнитной теоремы. Такое чтение дополнительно
предполагает точечное одинаковое парное ядро, отношение $\mu=\lambda$ и
отсутствие пространственных, дипольных и магнитных членов.
\end{nogocriterion}

\section{Что действительно известно о знаке}

Для протонного рисунка $uud$ и нейтронного $udd$ имеем
\begin{equation}
 E_p=\frac{\mu}{T},
 \qquad
 E_n=\frac{2(\mu-\lambda)}{3T},
 \qquad
 E_n-E_p=-\frac{\mu+2\lambda}{3T}.
 \label{eq:v8-baryon-em-neutron-proton-sign}
\end{equation}
Следовательно, знак отрицателен во всём положительном электростатическом
конусе $\mu>0$, $\lambda>0$. Это более устойчивое условное утверждение,
чем специальная точка $\mu=\lambda=1$. Однако дополнительные магнитные и
пространственные члены способны изменить итог, поэтому физический знак
полной разности масс из данного тождества не следует.

\begin{protocol}
\begin{enumerate}
 \item арифметика: \textbf{точная, на всех $216$ состояниях};
 \item $A+C=Q_{\rm tot}^2$: \textbf{тождество};
 \item эпсилон-рисунок $(4,1,0,1)$: \textbf{точно};
 \item следовая норма $T$: \textbf{выведена точно};
 \item абсолютная электромагнитная энергия: \textbf{не выведена};
 \item условие свёртки к полному заряду: \textbf{$\mu=\lambda=1$};
 \item пространственные и магнитные члены: \textbf{не включены};
 \item знак при $\mu,\lambda>0$: \textbf{$E_n-E_p<0$};
 \item безусловная физическая теорема о массах: \textbf{не получена};
 \item следующий гейт: \textbf{происхождение пространственного ядра}.
\end{enumerate}
\end{protocol}

Машинный сертификат сохранён в
\path{s2t_v8_baryon_em_total_charge_identity_scope_gate_results.json}.